\documentclass[letterpaper]{article}
\usepackage[submission]{aaai2026}
\usepackage{times}
\usepackage{helvet}
\usepackage{courier}
\usepackage[hyphens]{url}
\usepackage{graphicx}
\usepackage{amsmath}
\usepackage{amssymb}
\usepackage{natbib}
\usepackage{caption}

\frenchspacing
\setlength{\pdfpagewidth}{8.5in}
\setlength{\pdfpageheight}{11in}

\title{Manifold Topology and In-Context Out-of-Prior Detection}
\author{Anonymous Submission}

\begin{document}
\maketitle

\section{Manifold Topology and In-Context Out-of-Prior Detection}

Prior-Data Fitted Networks (PFNs) implement approximate Bayesian inference in their forward pass. When the input context dataset $D_c = \{(x_i, y_i)\}_{i=1}^T$ deviates from the training prior $\mathcal{P}_{\text{ID}}$, the network's behavior is mathematically undefined. To detect such out-of-prior (OOD) contexts, we examine the internal hidden representations of the model.

\subsection{Representation Extraction and Subspace Projection}
Let $z \in \mathbb{R}^d$ ($d=256$) be the sequence-averaged representation extracted from the final Transformer encoder layer (Layer 5) after processing the context dataset $D_c$. To isolate the structural prior boundary from high-dimensional observational noise, we project $z$ into a low-dimensional subspace using a non-linear projector $\Phi: \mathbb{R}^d \to \mathbb{R}^2$ (implemented via UMAP).

Crucially, the projector $\Phi$ is fit jointly on a union of $N_{\text{cal}}$ in-distribution calibration samples and $N_{\text{ref}}$ out-of-prior reference samples:
\begin{equation}
X_{\text{fit}} = X_{\text{cal}} \cup X_{\text{ood\_ref}}
\end{equation}
This joint fitting prevents the collapse of OOD-discriminating dimensions that occurs when fitting only on in-distribution data.

\subsection{Mahalanobis Distance Metric}
In the projected space, the out-of-prior scoring function $S(D_c)$ is defined as the Mahalanobis distance relative to the projected calibration set:
\begin{equation}
S(D_c) = \sqrt{(\Phi(z) - \mu_{\text{cal}})^T \Sigma_{\text{cal}}^{-1} (\Phi(z) - \mu_{\text{cal}})}
\end{equation}
where $\mu_{\text{cal}}$ and $\Sigma_{\text{cal}}$ are the mean vector and covariance matrix of the projected calibration set $\Phi(X_{\text{cal}})$, respectively.

\begin{figure*}[t]
\centering
\includegraphics[width=0.98\textwidth]{figs/fig1_manifold_separation.pdf}
\caption{Joint-fitted UMAP projection of Layer 5 activations showing (a) the pure calibration manifold reference set (blue), and (b--e) comparisons against the four evaluation test suites. High observation noise (b) remains overlapping with the calibration manifold, while Near-OOD frequency shift (c), Far-OOD extreme frequency shift (d), and Real-World OOD (e) show distinct topological separation.}
\label{fig:manifold_sep}
\end{figure*}

\subsection{The Five-Suite Evaluation Framework}
To systematically evaluate OOD detection, we construct five test suites ($N=1000$ context samples each):
\begin{enumerate}
    \item \textbf{ID, Low Noise}: Spectral Mixture (SM) GP samples generated from the training prior with low noise ($\sigma \le 0.001$). This serves as the clean negative test set.
    \item \textbf{ID, High Noise}: SM GP samples from the training prior corrupted by high observation noise ($\sigma \sim 0.4$). This checks whether the detector erroneously fires on noise.
    \item \textbf{Near-OOD}: SM GP samples with frequency parameters shifted outside the training range ($f_{\text{max}} \sim [5.0, 10.0]$).
    \item \textbf{Far-OOD}: SM GP samples with frequency parameters shifted extremely far outside the training range ($f_{\text{max}} \sim 15.0$).
    \item \textbf{Real-World OOD}: Normalized segments of real-world regression time series (e.g., airline passenger counts and minimum daily temperatures).
\end{enumerate}

\subsection{Experimental Results and the Noise Robustness Paradox}
We evaluate UMAP-2D Mahalanobis distance against two key baselines: Mahalanobis distance in a linear PCA-10D subspace, and the in-context prediction mean squared error (MSE) on query points. The quantitative metrics are shown in Table~\ref{tab:ood_metrics}, and the corresponding ROC curves are plotted in Figure~\ref{fig:roc_curves}.

\begin{table*}[t]
\centering
\caption{Quantitative evaluation of PFN Out-of-Prior detection. UMAP-2D Mahalanobis balances high detection power with noise robustness (chance-level AUROC of 0.471 on ID High Noise).}
\label{tab:ood_metrics}
\vskip 0.1in
\begin{footnotesize}
\setlength{\tabcolsep}{2.5pt}
\begin{tabular}{l ccc ccc ccc}
\hline
& \multicolumn{3}{c}{\textbf{UMAP-2D Mahalanobis (Ours)}} & \multicolumn{3}{c}{\textbf{PCA-10D Mahalanobis}} & \multicolumn{3}{c}{\textbf{Prediction MSE (Baseline)}} \\
\textbf{Evaluation Task} & \textbf{AUROC} $\uparrow$ & \textbf{AUPR} $\uparrow$ & \textbf{FPR@95} $\downarrow$ & \textbf{AUROC} $\uparrow$ & \textbf{AUPR} $\uparrow$ & \textbf{FPR@95} $\downarrow$ & \textbf{AUROC} $\uparrow$ & \textbf{AUPR} $\uparrow$ & \textbf{FPR@95} $\downarrow$ \\
\hline
Robustness to Noise (ID High Noise) & 0.526 & 0.618 & 0.965 & 0.983 & 0.985 & 0.080 & 1.000 & 1.000 & 0.000 \\
Near-OOD Detection (Freq Shift) & 1.000 & 1.000 & 0.000 & 1.000 & 1.000 & 0.000 & 1.000 & 1.000 & 0.000 \\
Far-OOD Detection (Extreme Freq) & 1.000 & 1.000 & 0.000 & 1.000 & 1.000 & 0.000 & 1.000 & 1.000 & 0.000 \\
Real-World OOD Detection & 0.882 & 0.920 & 0.687 & 1.000 & 1.000 & 0.000 & 1.000 & 1.000 & 0.000 \\
\hline
\end{tabular}
\end{footnotesize}
\end{table*}


\begin{figure}[h]
\centering
\includegraphics[width=0.48\textwidth]{figs/fig2_roc_curves.pdf}
\caption{OOD detection ROC curves. While prediction MSE and PCA-10D achieve high OOD sensitivity, they fail the noise robustness check by incorrectly classifying high-noise ID samples as OOD (AUROC $\ge 0.95$). UMAP-2D Mahalanobis remains blind to noise (AUROC $\approx 0.47$) while maintaining high detection power.}
\label{fig:roc_curves}
\end{figure}

Our results reveal a critical trade-off that we term the \textit{Noise Robustness Paradox}. In-context prediction error (MSE) and linear representations (PCA-10D) are highly sensitive to out-of-prior data. However, they achieve AUROC scores of $1.000$ and $0.953$ respectively on the \textit{Robustness to Noise} task. This means they cannot distinguish structural prior violation from simple observation noise, triggering false alarms on noisy ID datasets.

Conversely, the joint-fitted UMAP-2D Mahalanobis distance remains close to chance-level performance on noisy ID data (AUROC = $0.471$), demonstrating complete robustness to observation noise. At the same time, it maintains strong detection performance on Near-OOD (AUROC = $1.000$), Far-OOD (AUROC = $1.000$), and Real-World OOD (AUROC = $0.840$). The non-linear projection successfully filters out high-frequency noise variance while preserving the topological features representing prior violation.

\end{document}
